% !TeX encoding = UTF-8
% Français de France. Charger un seul archétype depuis le document principal.
\newcommand{\cvsummary}{Ingénieur mécanicien formé à NYU Abu Dhabi, préparant un master en mécatronique, avec une expérience de terrain en systèmes électriques industriels, automatisation et numérisation. Expérience en diagnostic de pannes, amélioration de l'efficacité opérationnelle et développement d'outils numériques pour les processus d'ingénierie.}

% --- COMPÉTENCES ---
\newcommand{\skills}{
    \section*{COMPÉTENCES TECHNIQUES}
    {\small
    \textbf{Commande et automatisation :} Commande PID avec anti-windup et gestion de la saturation des actionneurs, bases de représentation d'état et de commande LQR, machines à états finis, programmation d'automates, commande de moteurs à courant continu, commande embarquée en temps réel, MATLAB/Simulink.

    \vspace{1pt}
    \textbf{Systèmes embarqués et électronique :} Microcontrôleurs (Arduino/C++), FPGA, PWM, ADC/DAC, codeurs en quadrature, pilotes de moteurs, IMU et fusion de capteurs, Raspberry Pi, IoT, capteurs et instrumentation, conception logique numérique (Quartus), intégration matériel-logiciel.

    \vspace{1pt}
    \textbf{Systèmes électriques :} Schémas unifilaires, analyse de postes haute tension (harmoniques), compensation du facteur de puissance (batteries de condensateurs), AutoCAD Electrical, systèmes industriels de gestion d'entrepôt (WMS) et de traçabilité (Industrie 4.0).

    \vspace{1pt}
    \textbf{Conception mécanique et simulation :} SolidWorks (CSWA), Solid Edge, CATIA V5, Siemens NX, CREO, ANSYS (Mechanical, Fluent), analyse par éléments finis (FEA), mécanique des fluides numérique (CFD), conception de systèmes thermiques (TEC, CVC).

    \vspace{1pt}
    \textbf{Fabrication et qualité :} CNC, fabrication additive et FDM, injection plastique (analyse Moldflow), matériaux composites, cotation dimensionnelle et géométrique (GD\&T), contrôle qualité ISO 9001, Lean Six Sigma White Belt.

    \vspace{1pt}
    \textbf{Programmation :} C++, Python, MATLAB, communication Qt/TCP avec des contrôleurs de moteurs, JavaScript, SQL, Git.

    \vspace{1pt}
    \textbf{Gestion de projets :} MS Project, Jira, Notion, coordination pluridisciplinaire, rapports techniques.

    \vspace{1pt}
    \textbf{Langues :} Espagnol langue maternelle, anglais C2, français C2, italien A1, allemand A1.
    }
}   

% --- PROJETS ---
\newcommand{\vawttitle}{Éolienne à axe vertical : récupération d'énergie éolienne urbaine}
\newcommand{\vawtdates}{09/2024 - 05/2025}
\newcommand{\vawtlocation}{Abou Dabi, ÉAU}
\newcommand{\vawtdescOne}{Développement d'une nouvelle éolienne à axe vertical à pales combinant traînée et portance, adaptée au vent urbain des ÉAU, par conception mécanique et simulations CFD pour l'optimisation énergétique. Brevet en cours.}
\newcommand{\vawtdescTwo}{Évaluation de l'aérodynamique du rotor par modélisation du système de récupération d'énergie éolienne et analyse des résultats de simulations CFD sous ANSYS Fluent.}

\newcommand{\fsaetitle}{Formula SAE -- Suspension et châssis en fibre de carbone}
\newcommand{\fsaedates}{Automne 2024 -- Aujourd'hui}
\newcommand{\fsaelocation}{Abou Dabi, ÉAU}
\newcommand{\fsaedescOne}{Développement des systèmes porteurs du véhicule, dont la suspension et le châssis en fibre de carbone, par modélisation sous SolidWorks et analyse par éléments finis sous ANSYS.}
\newcommand{\fsaedescTwo}{Évaluation du mouvement de la suspension par simulation du mécanisme et analyses cinématiques sous SolidWorks et ANSYS.}

\newcommand{\tactitle}{Mini-réfrigérateur -- Système thermoélectrique TEC à deux étages}
\newcommand{\tacdates}{Printemps 2025}
\newcommand{\taclocation}{Abou Dabi, ÉAU}
\newcommand{\tacdescOne}{Développement d'un dispositif de refroidissement thermoélectrique à deux étages sans fluide frigorigène, en remplaçant le circuit frigorifique par des modules TEC.}
\newcommand{\tacdescTwo}{Intégration de capteurs et d'une instrumentation en temps réel dans l'ensemble de refroidissement pour suivre son comportement thermique en fonctionnement.}

\newcommand{\orbitaltitle}{OrbitalGrid -- Système SIG pour applications spatiales}
\newcommand{\orbitaldates}{06/2024 - Aujourd'hui}
\newcommand{\orbitallocation}{Abou Dabi, ÉAU}
\newcommand{\orbitaldescOne}{Développement d'un système d'intégration de données spatiales à travers un projet de cartographie SIG fondé sur SQL, combinant numérisation 3D et données satellitaires.}
\newcommand{\orbitaldescTwo}{Mise en place d'une gestion de données géospatiales sous PostgreSQL, en structurant les informations pour la cartographie aérospatiale.}

\newcommand{\softrobtitle}{Exosquelette robotique pour la rééducation après un AVC}
\newcommand{\softrobdates}{04/2024 - 05/2024}
\newcommand{\softroblocation}{Abou Dabi, ÉAU}
\newcommand{\softrobdescOne}{Conception d'un mécanisme de rééducation actionné, sous la forme d'un exosquelette de doigts inspiré de l'origami intégrant des fils en alliage à mémoire de forme dans une structure souple.}
\newcommand{\softrobdescTwo}{Intégration de matériaux souples et d'éléments actifs dans une conception associant polymère Dragon Skin et fils SMA, combinant interfaces déformables et actionnement thermosensible.}

\newcommand{\exoplanettitle}{Simulateur de transits d'exoplanètes}
\newcommand{\exoplanetdates}{Printemps 2025}
\newcommand{\exoplanetlocation}{Abou Dabi, ÉAU}
\newcommand{\exoplanetdescOne}{Développement d'un modèle physique paramétrable de transits d'exoplanètes en Python, par implémentation de la mécanique orbitale et de la physique stellaire.}
\newcommand{\exoplanetdescTwo}{Génération de données d'observation synthétiques sous forme de séries temporelles de courbes de lumière, par calcul des transits à partir de paramètres astrophysiques.}

\newcommand{\windTurbinetitle}{Éolienne flottante en mer -- Stabilisation active}
\newcommand{\windTurbinedates}{05/2026 - 06/2026}
\newcommand{\windTurbinelocation}{Besançon, France}
\newcommand{\windTurbinedescOne}{Conception et réalisation sous SolidWorks d'un prototype d'éolienne semi-submersible à trois colonnes à l'échelle 1:1000, avec stabilité en tangage et roulis et GM = 5,38 cm ; intégration d'un MPU6050 et d'un cardan actionné par servomoteurs pour stabiliser activement la nacelle lors d'essais en bassin à houle.}
\newcommand{\windTurbinedescTwo}{Établissement de la stabilité passive en tangage, avec GM = 5,38 cm et raideur de rappel K = 0,0607 N$\cdot$m/rad, par calcul hydrostatique d'un prototype semi-submersible à trois colonnes à l'échelle 1:1000.}

\newcommand{\carbonmetertitle}{CARBONMETRE -- Écoconception et mesure carbone}
\newcommand{\carbonmeterdates}{09/2025 - 12/2025}
\newcommand{\carbonmeterlocation}{Besançon, France}
\newcommand{\carbonmeterdescOne}{Application de principes métrologiques à la mesure carbone dans le projet d'ingénierie CARBONMETRE, en combinant mesure, fabrication et écoconception.}
\newcommand{\carbonmeterdescTwo}{Mise en relation de la fabrication et de l'ingénierie environnementale par application des principes de fabrication industrielle à un projet de mesure carbone.}

\newcommand{\hvactitle}{Amélioration thermique d'une chambre en résidence -- Refroidissement passif et évaporatif}
\newcommand{\hvacdates}{06/2026}
\newcommand{\hvaclocation}{Besançon, France}
\newcommand{\hvacdescOne}{Réduction déclarée de la température d'une chambre de 9 m$^2$ jusqu'à 3 \textdegree{}C par rapport à la situation avant intervention, grâce à un écran solaire extérieur et à un refroidisseur à eau gelée entraîné par ventilateur.}
\newcommand{\hvacdescTwo}{Blocage de l'exposition solaire aux heures de pointe, de 14 h à 21 h, par superposition à l'extérieur de papier aluminium réfléchissant, de PSE de 1 cm et de carton ondulé.}

\newcommand{\controlBoxtitle}{Installation électrique complète -- Automate programmable pour événements de taille moyenne}
\newcommand{\controlBoxdates}{06/2023 - 08/2023}
\newcommand{\controlBoxlocation}{Mexico, Mexique}
\newcommand{\controlBoxdescOne}{Réduction de 30\% du temps de montage du système électrique par rapport à l'installation précédente, grâce à la conception d'un coffret de commande à automate programmable de 12 entrées pour installations événementielles.}
\newcommand{\controlBoxdescTwo}{Développement d'une commande électrique à automate programmable pour une installation événementielle de taille moyenne, adaptée aux exigences de fonctionnement des équipements audiovisuels.}

\newcommand{\hyperlooptitle}{NYC Hyperloop -- Véhicule à échelle réduite}
\newcommand{\hyperloopdates}{01/2024 - 05/2024}
\newcommand{\hyperlooplocation}{New York, États-Unis}
\newcommand{\hyperloopdescOne}{Pilotage du développement mécanique en tant que responsable du sous-groupe mécanique d'un projet Hyperloop à échelle réduite, pour un concept de liaison New York--Long Island.}
\newcommand{\hyperloopdescTwo}{Évaluation de la lévitation magnétique par modélisation multiphysique du système de transport et simulations sous COMSOL et ANSYS.}

\newcommand{\dronetitle}{Quadricoptère fonctionnel -- Structure allégée}
\newcommand{\dronedates}{03/2024 - 05/2024}
\newcommand{\dronelocation}{New York, États-Unis}
\newcommand{\dronedescOne}{Réduction de 40\% du poids du quadricoptère par rapport à la conception de référence du projet, grâce à une structure imprimée en 3D en PLA et TPU.}
\newcommand{\dronedescTwo}{Évaluation de la résistance aux chocs de la structure par essais de chute comparés aux critères de performance pendant le développement.}

\newcommand{\trentxwbtitle}{Rolls-Royce Trent XWB -- Modèle 3D fonctionnel}
\newcommand{\trentxwbdates}{01/2024 - 05/2024}
\newcommand{\trentxwblocation}{New York, États-Unis}
\newcommand{\trentxwbdescOne}{Codéveloppement d'un assemblage CAO intégré de propulsion, avec modèle 3D fonctionnel complet du Trent XWB, par modélisation des composants du moteur sous SolidWorks et Fusion 360.}
\newcommand{\trentxwbdescTwo}{Évaluation du comportement aérodynamique du modèle de moteur par analyses et simulations sous ANSYS.}

\newcommand{\liostitle}{LIOS -- Assistant vocal avec IA}
\newcommand{\liosdates}{12/2024 - Aujourd'hui}
\newcommand{\lioslocation}{Abou Dabi, ÉAU}
\newcommand{\liosdescOne}{Développement d'un assistant vocal avec IA en Python pour automatiser des tâches, en intégrant traitement automatique du langage, données en temps réel et API externes.}
\newcommand{\liosdescTwo}{Mise en place d'interactions tenant compte du contexte par conservation des informations pertinentes dans les processus de l'assistant au moyen d'une mémoire contextuelle.}

\newcommand{\fitpaltitle}{FitPal -- Dispositif intelligent d'entraînement physique}
\newcommand{\fitpaldates}{09/2021 - 12/2021}
\newcommand{\fitpallocation}{Abou Dabi, ÉAU}
\newcommand{\fitpaldescOne}{Codéveloppement d'un dispositif de mesure embarqué sur M5Stack pour l'analyse d'exercices en temps réel, intégrant acquisition et traitement dans un produit d'entraînement physique.}
\newcommand{\fitpaldescTwo}{Connexion du matériel à une interface mobile par intégration d'Arduino IoT App et de fonctions de transfert de données IoT.}

\newcommand{\covidtitle}{Cabine de désinfection automatisée}
\newcommand{\coviddates}{09/2020 - 02/2021}
\newcommand{\covidlocation}{Mexico, Mexique}
\newcommand{\coviddescOne}{Développement d'un prototype matériel de cabine de désinfection pilotée par Arduino, intégrant électronique de commande et système de brumisation.}
\newcommand{\coviddescTwo}{Ajout d'une communication sans fil au système par connexion sans câble des éléments de commande du prototype.}

\newcommand{\rovertitle}{Rover téléopéré -- Numérisation spatiale 3D}
\newcommand{\roverdates}{09/2026 - Aujourd'hui}
\newcommand{\roverlocation}{Oviedo, Espagne}
\newcommand{\roverdescOne}{Développement en cours d'une plateforme mobile de mesure téléopérée, avec intégration visée de la commande embarquée, de l'acquisition 3D et de la visualisation à distance dans un rover.}
\newcommand{\roverdescTwo}{Développement en cours de la commande embarquée du rover par intégration du sous-système de commande, avec le fonctionnement télécommandé comme jalon cible.}

\newcommand{\microsensortitle}{Microcapteur piézorésistif de pression de détonation}
\newcommand{\microsensordates}{05/2026 - 06/2026}
\newcommand{\microsensorlocation}{Besançon, France}
\newcommand{\microsensordescOne}{Conception d'un microcapteur de pression piézorésistif, rédaction d'instructions détaillées de fabrication couvrant la pulvérisation cathodique et le lift-off, et automatisation en Python des calculs géométriques avec contraintes de facteur de jauge, $\Delta R/R_0$ et résistivité.}
\newcommand{\microsensordescTwo}{Pilotage de la sélection de matériaux sous contraintes électromécaniques pour une conception répondant à une interface de 50 $\Omega$ et à une tenue de 10 GPa, en conciliant exigences électriques, mécaniques et économiques à l'échelle du wafer.}

\newcommand{\missiletitle}{Simulateur de dynamique d'interception de missiles}
\newcommand{\missilelocation}{Besançon, France}
\newcommand{\missiledescOne}{Développement sous MATLAB d'une simulation 3D d'interception de missiles avec loi de guidage par navigation proportionnelle, intégration numérique et visualisation des trajectoires en temps réel avec détection d'impact.}

% --- EXPÉRIENCE ---
\newcommand{\carbonorg}{Energy Innovation Capital}
\newcommand{\carbontitle}{Analyste de marché -- Captage du carbone}
\newcommand{\carbondates}{10/2023 - 01/2024}
\newcommand{\carbonlocation}{Abou Dabi, ÉAU}
\newcommand{\carbonbullets}{%
  \item Appui aux décisions d'investissement par recommandations stratégiques sur des jeunes entreprises CCUS, grâce à l'analyse de technologies de captage du carbone pour un fonds de capital-risque d'Abou Dabi.
  \item Traduction de recherches techniques à destination des décideurs par présentations aux analystes seniors, synthétisant les avancées scientifiques du captage du carbone.
  \item Clarification des moteurs technologiques et commerciaux par analyses des tendances du marché combinant recherche technique et données sectorielles.
  \item Contribution à l'évaluation du risque d'investissement du fonds par examen conjoint des évolutions technologiques CCUS et des conditions de marché.
}

\newcommand{\spacelaborg}{NYU Abu Dhabi, Centre des sciences spatiales et d'astrophysique}
\newcommand{\spacelabtitle}{Ingénieur et chercheur R\&D -- Plateforme de sciences planétaires}
\newcommand{\spacelabdates}{05/2025 - 09/2025}
\newcommand{\spacelablocation}{Abou Dabi, ÉAU}
\newcommand{\spacelabbullets}{%
  \item Amélioration de 40\% de l'efficacité du traitement d'images par rapport à la chaîne précédente, grâce à l'optimisation de la compression, de la reprojection cartographique et de la correction d'images assistée par apprentissage automatique.
  \item Transformation d'images de sonde spatiale en \href{https://romvio.github.io/mars-viewer/}{application web} interactive sous React et JavaScript, rendant les données de la mission HOPE accessibles par une interface cartographique visuelle ; adaptation de la plateforme à iOS et soumission à l'App Store d'Apple aux ÉAU.
  \item Assemblage d'images de plusieurs passages orbitaux en cartes planétaires cohérentes par alignement des recouvrements, correction des différences de luminosité et implémentation des projections cylindrique et Robinson.
}

\newcommand{\geoorg}{Global Education NYUAD}
\newcommand{\geotitle}{Chargé de coordination Global Education -- Étudiants en mobilité entrante}
\newcommand{\geodates}{12/2025 - 05/2025}
\newcommand{\geolocation}{Abou Dabi, ÉAU}
\newcommand{\geobullets}{%
  \item Coordination des opérations de programmes de terrain aux ÉAU, avec sorties universitaires et culturelles pour étudiants en mobilité entrante à NYU Abu Dhabi.
  \item Organisation de la réalisation des activités par encadrement de l'équipe de coordinateurs et planification logistique transversale.
  \item Contribution à la sécurité des déplacements de groupe par protocoles et précautions intégrés à la logistique des sorties.
  \item Accompagnement des participants en tant que référent étudiant, par coordination des échanges entre étudiants et services universitaires.
}

\newcommand{\geoOneorg}{Global Education NYUAD}
\newcommand{\geoOnetitle}{Chargé de coordination Global Education -- Étudiants de première année}
\newcommand{\geoOnedates}{09/2026 - 12/2024}
\newcommand{\geoOnelocation}{Abou Dabi, ÉAU}
\newcommand{\geoOnebullets}{%
  \item Accompagnement de l'intégration en première année par animation de séances d'accueil et de découverte de la vie universitaire.
  \item Conseil continu par mentorat et accompagnement entre pairs pour aider les étudiants à utiliser les ressources du campus.
  \item Transmission d'informations pratiques sur les installations et les services universitaires au cours de visites du campus.
  \item Contribution à l'intégration des promotions par planification d'événements universitaires et culturels inclusifs pour les étudiants de première année.
}

\newcommand{\hcorg}{Centre de santé de NYUAD}
\newcommand{\hctitle}{Emploi étudiant -- Responsable des réseaux sociaux}
\newcommand{\hcdates}{09/2024 - 05/2025}
\newcommand{\hclocation}{Abou Dabi, ÉAU}
\newcommand{\hcbullets}{%
  \item Mobilisation des étudiants par campagnes ciblées de bien-être et adaptation de la communication de santé au public de NYU Abu Dhabi.
  \item Création de supports de communication numérique et de contenus publiés sur les réseaux sociaux pour promouvoir les initiatives de santé.
  \item Rôle de contact principal pour les demandes étudiantes, avec coordination des réponses aux questions adressées au centre de santé.
  \item Contribution à la résolution des demandes en associant accompagnement direct des étudiants et communication du centre de santé.
}

\newcommand{\alpesorg}{Alpes Technologies México}
\newcommand{\alpestitle}{Ingénieur projets junior -- Gestion d'entrepôt et systèmes électriques}
\newcommand{\alpesdates}{06/2024 - 09/2024}
\newcommand{\alpeslocation}{Mexico, Mexique}
\newcommand{\alpesbullets}{%
  \item Fabrication et installation de batteries de condensateurs industrielles avec filtres harmoniques après analyse des harmoniques et pics de tension des postes, atteignant un facteur de puissance de 0,99 et réduisant les coûts énergétiques de 15\% et les arrêts électriques de 25\%.
  \item Remplacement du suivi manuel des stocks par un système de gestion d'entrepôt avec codes QR, développé en HTML/JavaScript et Google Drive, améliorant de 60\% l'efficacité opérationnelle et de fabrication.
  \item Analyse des dommages électriques après l'incendie d'un complexe hôtelier par mesures sur site, avec présentation de recommandations correctives à la direction générale.
}

\newcommand{\epigsaorg}{EPIGSA}
\newcommand{\epigsatitle}{Assistant en ingénierie des données -- SIG et hydraulique}
\newcommand{\epigsadates}{05/2024 - 07/2024}
\newcommand{\epigsalocation}{Mexico, Mexique}
\newcommand{\epigsabullets}{%
  \item Réduction de 25\% des erreurs dans les données SIG par automatisation des contrôles qualité pour des projets municipaux d'infrastructures de traitement des eaux.
  \item Maintenance et validation des données d'infrastructures d'ingénierie sous PostgreSQL, assurant cohérence et accessibilité pour les équipes pluridisciplinaires.
}

\newcommand{\solarshiporg}{Solar Ship}
\newcommand{\solarshiptitle}{Stagiaire en stratégie et conception de mobilité aérienne}
\newcommand{\solarshipdates}{05/2023 - 08/2023}
\newcommand{\solarshiplocation}{Toronto, Canada}
\newcommand{\solarshipbullets}{%
  \item Développement de concepts de logistique aérienne pour des zones isolées de Zambie et de Tanzanie par planification d'infrastructures destinées à des communautés hors réseau.
  \item Contribution à la conception d'un système aéronautique entièrement électrique dans l'initiative Peace+Freedom pour améliorer la connectivité des communautés.
  \item Mise en relation de la conception du véhicule et des besoins d'infrastructure dans un concept de logistique aérienne isolée, en considérant les besoins d'accès avec les systèmes de l'aéronef.
  \item Intégration d'objectifs techniques et sociaux dans un projet multinational de mobilité, combinant ingénierie aérospatiale et exigences d'accès des communautés.
}

\newcommand{\crousorg}{CROUS Besançon Bourgogne Franche-Comté}
\newcommand{\croustitle}{Emploi étudiant -- Restauration universitaire}
\newcommand{\crousdates}{10/2024 - Aujourd'hui}
\newcommand{\crouslocation}{Besançon, France}
\newcommand{\crousbullets}{%
  \item Maintien de la continuité du travail en cuisine aux heures de pointe par appui à la préparation de repas universitaires en grand volume.
  \item Respect des procédures de nettoyage et prise en charge de tâches d'hygiène et de sécurité alimentaire en cuisine.
  \item Contribution à la continuité du service par réapprovisionnement et assistance pendant l'activité.
  \item Participation coordonnée aux tâches de cuisine et de service, avec adaptation aux besoins opérationnels soumis à des délais courts.
}

\newcommand{\rallyorg}{Championnat du monde des rallyes -- Guanajuato}
\newcommand{\rallytitle}{Commissaire technique officiel FIA}
\newcommand{\rallydates}{11/2019 - 04/2020}
\newcommand{\rallylocation}{Guanajuato, Mexique}
\newcommand{\rallybullets}{%
  \item Réalisation de contrôles techniques réglementaires et du marquage officiel des pneumatiques lors d'une épreuve du Championnat du monde des rallyes FIA, selon les procédures prescrites.
  \item Contribution à la conformité des véhicules de compétition par vérification des exigences de l'épreuve.
  \item Coordination des inspections avec les commissaires responsables et les pilotes sous la pression opérationnelle d'une compétition en cours.
  \item Application de procédures de marquage et de vérification dans des travaux de contrôle décrits conformément aux normes ISO 9001.
}

\newcommand{\papaloteorg}{Papalote Museo del Niño}
\newcommand{\papalotetitle}{Chargé d'accueil et de médiation}
\newcommand{\papalotedates}{05/2019 - 07/2019}
\newcommand{\papalotelocation}{Mexico, Mexique}
\newcommand{\papalotebullets}{%
  \item Présentation de systèmes techniques à des non-spécialistes lors de visites d'expositions scientifiques et technologiques, avec explications adaptées aux enfants et aux adultes.
  \item Proposition d'améliorations concrètes de l'expérience de visite, mises en œuvre pour le fonctionnement du musée.
  \item Accompagnement de visiteurs de différents âges en adaptant la communication à leurs besoins.
  \item Mise en relation de la science et de la participation pratique par démonstrations et explications accessibles de concepts techniques dans les expositions.
}
